\documentclass[11pt]{article}
\usepackage[a4paper,margin=25mm]{geometry}
\usepackage{amsmath,amssymb,amsthm}
\newtheorem{proposition}{Proposition}
\title{Exclusions using the marked generators of a regular cover}
\author{Human Compression project}\date{}
\begin{document}\maketitle
We acquire the normalizer description of chiral quotients in Hartley,
Hubard and Leemans, \emph{Two atlases of abstract chiral polytopes for
small groups}, Theorem~3.1. We keep every possible flag stabilizer,
including those with nontrivial core. The cover action is not assumed
faithful. Enumeration of normal subgroups remains an explicit software
premise. We also use a necessary condition already implicit in the
reconstruction of faces from a flag graph.

\begin{proposition}
Let a regular polytope with group $W=\langle r_0,\ldots,r_{d-1}\rangle$
cover a chiral abstract polytope, and let $N$ stabilize its base flag. Put
\[
 H=\langle r_0r_1,\ldots,r_0r_{d-1}\rangle,
 \qquad W_i=\langle r_j:j\ne i\rangle.
\]
Then $[W:H]=2$, $N_W(N)=H$, and
\[
 N\trianglelefteq H,\qquad N\ne1,\qquad N\not\trianglelefteq W.
\]
The kernel of the cover action is $N\cap N^{r_0}$, which may be
nontrivial. Moreover $\bigcap_{i=0}^{d-1}NW_i=N$.
\end{proposition}
\begin{proof}
The subgroup generated by the displayed pairs contains every product
of two generators, since
\[
 r_ir_j=(r_0r_i)^{-1}(r_0r_j),\qquad r_0r_0=1.
\]
It consists of the
even words and has index at most two. In a chiral flag graph every
adjacency interchanges the two automorphism orbits. An odd word cannot
stabilize a flag. The even subgroup therefore has index two and agrees
with the normalizer of $N$, by the acquired normalizer theorem.

The normal-subgroup conditions follow from the same theorem. Since
$N$ is normal in the subgroup $H$ of index two, its only conjugates in
$W$ are $N,N^{r_0}$. Their intersection is therefore the core and the
kernel of the coset action. If $N=1$, its normalizer would be all of $W$,
contrary to chirality.

Identify graph vertices with right cosets $Nw$, with colour $j$ sending
$Nw$ to $Nwr_j$. The face of rank $i$ through the base flag is precisely
the connected component obtained by omitting colour $i$. Its vertices
are the cosets $Nh$ with $h\in W_i$. Hence $Nw$ has every base face if
and only if $w\in NW_i$ for every $i$. A flag is determined by its faces;
such a vertex must equal $N$, so $w\in N$. The reverse inclusion is
immediate.
\end{proof}

For each literal source marking, the reader checks the generated group,
its order, the involutions, adjacent orders and
string commutations. The database identifier is a source label; no abstract
identification is used by this test. If $H=W$, this marking is excluded. Otherwise it
enumerates all normal subgroups of $H$ and retains exactly the candidates
satisfying the displayed group conditions. For each it computes all
$NW_i$. If $T_i$ is a right transversal of $N\cap W_i$ in $W_i$, then
$NW_i=\coprod_{t\in T_i}Nt$; this is the finite computation used.
A negative result also retains a literal element $w\notin N$ and checks
$n_i^{-1}w\in W_i$ for some $n_i\in N$ at every rank.

Passing this test does not prove that the quotient is a polytope.
Survivors still require the full face and flag axioms and comparison
up to coloured graph isomorphism. No new exclusion is credited before
the individual source entry has finished and been reviewed. The earlier filter
for faithful actions alone does not exclude arbitrary cover actions.
The source examples
and original normalizer enumeration retain their original attribution.
This component supplies an explicit finite application of the acquired
group and face criteria, not a new complete atlas theorem.
\end{document}
